Normalmente referido como \"casting\", as conversões de tipos de dados podem
ser feitas com as mais diferentes finalidades.

\subsubsection{Conversões implícitas}

Ao avaliar uma expressão, o compilador acaba por realizar um processo de conversão 
de dados que respeita uma regra chamada \textbf{promoção automática} onde todas as 
variáveis são convertidas para um mesmo tipo.

Essa regra funciona de forma que os tipos de menor tamanho são convertidos para os 
tipos de maior tamanho que existem dentro da expressão, essa hierarquia, por assim 
dizer, de tipos segue a seguinte ordem:

\begin{verbatim}
    char < short < int < long < long long < float < double < long double
\end{verbatim}

A esse ponto você deve estar se perguntando o porquê de existirem todos esses tipos 
se no final eles são todos interpretados como iguais pelo compilador.

\begin{ps}
    Não entenda o processo de promoção automática como se todas as variáveis fossem 
    tratadas como iguais pelo computador e o processo de tipagem como apenas uma 
    regalia para facilitar o processo de desenvolvimento.
\\\\
    A tipagem de variáveis continua sendo extremamente necessária para que o 
    compilador realize, de maneira correta e precisa, processos que eu provavelmente 
    sequer me atreveria a tentar explicar neste livro.
\\\\
    Sendo assim, o processo de promoção automática faz com que as variáveis \"sejam 
    tratadas como iguais\" somente dentro do escopo de uma expressão ou operação, regra essa que 
    não se aplica ao programa em um escopo maior.
\end{ps}
